% GENERATED FILE. Edit canonical lessons, then run npm run generate:books.
% canonical-source-hash: fnv1a64:38e9abbca02767ad
% canonical-lessons: BN-C445-shok, BN-C445-kanna, BN-C445-dirghoshvas, BN-C445-phopano, BN-C445-bilap-kora, BN-R445-review-438-445

\chapter{Grief and Tears}
\label{ch:bn-grief-and-tears}
\hlchaptermodality{445}

\begin{chapteropening}
\textbf{By the end of this chapter:} \emph{I can talk about grief, tears and sighing.}

Complete the last lesson of chapter 445: I can talk about grief, tears and sighing.
\end{chapteropening}

\section[\texorpdfstring{\bn{শোক} (shok)}{শোক (shok)}]{\bn{শোক} (shok) — grief; mourning}

\label{lesson:BN-C445-shok}



\begin{quote}
Before the new one: say the Bengali for to forgive, then the Bengali for ill feeling.
\end{quote}

\subsection*{You'll want to know: \bn{শোক}}
\textbf{\bn{শোক}} — \emph{shok} — \textquotedblleft{}grief; mourning\textquotedblright{}.

\begin{quote}
\bn{দাদুর} \bn{মৃত্যুর} \bn{পরে} \bn{পুরো} \bn{পরিবার} \bn{শোকে} \bn{ছিল।}
\end{quote}

\begin{quote}
\emph{dādur mrityur pôre puro pôribār shoke chhilô.}
\end{quote}

\begin{quote}
After my grandfather's death the whole family was in mourning.
\end{quote}

\begin{cousinweb}
\textbf{\bn{শোক}} (\emph{shok}) is a tatsama, Sanskrit \emph{śoka}, \textquotedblleft{}burning\textquotedblright{}, and so \textquotedblleft{}sorrow, grief\textquotedblright{} (Monier-Williams).
\end{cousinweb}

\begin{grammarlens}[title={What you've built}]
One more word for describing what you lived and felt.
\end{grammarlens}

\subsection*{Your turn}
\begin{itemize}
\raggedright
  \item \emph{Say it:} \emph{shok}
  \item \emph{Say it:} \emph{shok}, once more
  \item \emph{Say it:} \emph{dādur mrityur pôre puro pôribār shoke chhilô.} — After my grandfather's death the whole family was in mourning.
  \item \emph{Recall:} say the Bengali for to forgive, then the Bengali for ill feeling, then say \emph{shok} again
\end{itemize}

\begin{tcolorbox}[breakable,colback=teal!4,colframe=teal!35!black,title={Before you move on}]
What does \bn{শোক} mean? (\textquotedblleft{}Grief; mourning\textquotedblright{}.) Say it once more.
\end{tcolorbox}
\section[\texorpdfstring{\bn{কান্না} (kānnā)}{কান্না (kānnā)}]{\bn{কান্না} (kānnā) — crying; weeping}

\label{lesson:BN-C445-kanna}



\begin{quote}
Before the new one: say the Bengali for ill feeling, then the Bengali for grief.
\end{quote}

\subsection*{You'll want to know: \bn{কান্না}}
\textbf{\bn{কান্না}} — \emph{kānnā} — \textquotedblleft{}crying; weeping\textquotedblright{}.

\begin{quote}
\bn{রাতে} \bn{শিশুর} \bn{কান্না} \bn{থামল} \bn{না।}
\end{quote}

\begin{quote}
\emph{rāte shishur kānnā thāmlô nā.}
\end{quote}

\begin{quote}
In the night the baby's crying would not stop.
\end{quote}

\bn{কান্না} \bn{পাওয়া}, to feel like crying: \bn{আমার} \bn{কান্না} \bn{পেল}.

\begin{cousinweb}
\textbf{\bn{কান্না}} (\emph{kānnā}) is listed by CDIAL for Bengali, as lamentation, under Sanskrit \emph{kránda}, \textquotedblleft{}crying out\textquotedblright{} (CDIAL 3572); \bn{কাঁদা} (taught), to cry, is from the same root.
\end{cousinweb}

\begin{grammarlens}[title={What you've built}]
One more word for describing what you lived and felt.
\end{grammarlens}

\subsection*{Your turn}
\begin{itemize}
\raggedright
  \item \emph{Say it:} \emph{kānnā}
  \item \emph{Say it:} \emph{kānnā}, once more
  \item \emph{Say it:} \emph{rāte shishur kānnā thāmlô nā.} — In the night the baby's crying would not stop.
  \item \emph{Recall:} say the Bengali for ill feeling, then the Bengali for grief, then say \emph{kānnā} again
\end{itemize}

\begin{tcolorbox}[breakable,colback=teal!4,colframe=teal!35!black,title={Before you move on}]
What does \bn{কান্না} mean? (\textquotedblleft{}Crying; weeping\textquotedblright{}.) Say it once more.
\end{tcolorbox}
\section[\texorpdfstring{\bn{দীর্ঘশ্বাস} (dīrghôshvās)}{দীর্ঘশ্বাস (dīrghôshvās)}]{\bn{দীর্ঘশ্বাস} (dīrghôshvās) — a sigh; a long breath}

\label{lesson:BN-C445-dirghoshvas}



\begin{quote}
Before the new one: say the Bengali for grief, then the Bengali for crying.
\end{quote}

\subsection*{You'll want to know: \bn{দীর্ঘশ্বাস}}
\textbf{\bn{দীর্ঘশ্বাস}} — \emph{dīrghôshvās} — \textquotedblleft{}a sigh; a long breath\textquotedblright{}.

\begin{quote}
\bn{পুরনো} \bn{ছবিটা} \bn{দেখে} \bn{দাদু} \bn{একটা} \bn{দীর্ঘশ্বাস} \bn{ফেললেন।}
\end{quote}

\begin{quote}
\emph{purôno chhôbiṭā dekhe dādu ekṭā dīrghôshvās phellen.}
\end{quote}

\begin{quote}
Looking at the old photo, my grandfather heaved a sigh.
\end{quote}

\bn{দীর্ঘশ্বাস} \bn{ফেলা}, to sigh.

\begin{cousinweb}
\textbf{\bn{দীর্ঘশ্বাস}} (\emph{dīrghôshvās}) is Sanskrit \emph{dīrgha}, \textquotedblleft{}long\textquotedblright{}, and \emph{śvāsa}, \textquotedblleft{}breath\textquotedblright{} (Monier-Williams): a long breath. The second half is the \bn{শ্বাস} inside \bn{নিঃশ্বাস} (taught), breath.
\end{cousinweb}

\begin{grammarlens}[title={What you've built}]
One more word for describing what you lived and felt.
\end{grammarlens}

\subsection*{Your turn}
\begin{itemize}
\raggedright
  \item \emph{Say it:} \emph{dīrghôshvās}
  \item \emph{Say it:} \emph{dīrghôshvās}, once more
  \item \emph{Say it:} \emph{purôno chhôbiṭā dekhe dādu ekṭā dīrghôshvās phellen.} — Looking at the old photo, my grandfather heaved a sigh.
  \item \emph{Recall:} say the Bengali for grief, then the Bengali for crying, then say \emph{dīrghôshvās} again
\end{itemize}

\begin{tcolorbox}[breakable,colback=teal!4,colframe=teal!35!black,title={Before you move on}]
What does \bn{দীর্ঘশ্বাস} mean? (\textquotedblleft{}A sigh; a long breath\textquotedblright{}.) Say it once more.
\end{tcolorbox}
\section[\texorpdfstring{\bn{ফোঁপানো} (phõpāno)}{ফোঁপানো (phõpāno)}]{\bn{ফোঁপানো} (phõpāno) — to sob}

\label{lesson:BN-C445-phopano}



\begin{quote}
Before the new one: say the Bengali for crying, then the Bengali for a sigh.
\end{quote}

\subsection*{You'll want to know: \bn{ফোঁপানো}}
\textbf{\bn{ফোঁপানো}} — \emph{phõpāno} — \textquotedblleft{}to sob\textquotedblright{}.

\begin{quote}
\bn{বাবা} \bn{বকার} \bn{পরে} \bn{ছেলেটা} \bn{বিছানায়} \bn{শুয়ে} \bn{ফোঁপাচ্ছিল।}
\end{quote}

\begin{quote}
\emph{bābā bôkār pôre chheleṭā bichhānāy shuye phõpāchchhilô.}
\end{quote}

\begin{quote}
After his father scolded him, the boy lay on the bed sobbing.
\end{quote}

Past (he/she …): \bn{সেই} \bn{শিশু} \bn{ফোঁপাল}, \emph{sei shishu phõpālô}.

\begin{cousinweb}
\textbf{\bn{ফোঁপানো}} (\emph{phõpāno}) is an imitative verb, its sound standing for the catch of breath in sobbing. No dictionary we use records an older form, so its history is not settled.
\end{cousinweb}

\begin{grammarlens}[title={What you've built}]
One more word for describing what you lived and felt.
\end{grammarlens}

\subsection*{Your turn}
\begin{itemize}
\raggedright
  \item \emph{Say it:} \emph{phõpāno}
  \item \emph{Say it:} \emph{phõpāno}, once more
  \item \emph{Say it:} \emph{bābā bôkār pôre chheleṭā bichhānāy shuye phõpāchchhilô.} — After his father scolded him, the boy lay on the bed sobbing.
  \item \emph{Recall:} say the Bengali for crying, then the Bengali for a sigh, then say \emph{phõpāno} again
\end{itemize}

\begin{tcolorbox}[breakable,colback=teal!4,colframe=teal!35!black,title={Before you move on}]
What does \bn{ফোঁপানো} mean? (\textquotedblleft{}To sob\textquotedblright{}.) Say it once more.
\end{tcolorbox}
\section[\texorpdfstring{\bn{বিলাপ} \bn{করা} (bilāp kôrā)}{বিলাপ করা (bilāp kôrā)}]{\bn{বিলাপ} \bn{করা} (bilāp kôrā) — to lament; to wail in grief}

\label{lesson:BN-C445-bilap-kora}



\begin{quote}
Before the new one: say the Bengali for a sigh, then the Bengali for to sob.
\end{quote}

\subsection*{You'll want to know: \bn{বিলাপ} \bn{করা}}
\textbf{\bn{বিলাপ} \bn{করা}} — \emph{bilāp kôrā} — \textquotedblleft{}to lament; to wail in grief\textquotedblright{}.

\begin{quote}
\bn{স্বামীর} \bn{মৃত্যুর} \bn{পরে} \bn{মহিলাটি} \bn{বিলাপ} \bn{করছিলেন।}
\end{quote}

\begin{quote}
\emph{svāmīr mrityur pôre môhilāṭi bilāp kôrchhilen.}
\end{quote}

\begin{quote}
After her husband's death the woman was wailing in grief.
\end{quote}

Past (he/she …): \bn{তিনি} \bn{সারারাত} \bn{বিলাপ} \bn{করলেন}, \emph{tini sārārāt bilāp kôrlen}.

\begin{cousinweb}
\textbf{\bn{বিলাপ} \bn{করা}} (\emph{bilāp kôrā}) is Sanskrit \emph{vilāpa}, \textquotedblleft{}lamentation\textquotedblright{} (Monier-Williams), with \emph{kôrā}, \textquotedblleft{}to do\textquotedblright{} (CDIAL 2814).
\end{cousinweb}

\begin{grammarlens}[title={What you've built}]
One more word for describing what you lived and felt.
\end{grammarlens}

\subsection*{Your turn}
\begin{itemize}
\raggedright
  \item \emph{Say it:} \emph{bilāp kôrā}
  \item \emph{Say it:} \emph{bilāp kôrā}, once more
  \item \emph{Say it:} \emph{svāmīr mrityur pôre môhilāṭi bilāp kôrchhilen.} — After her husband's death the woman was wailing in grief.
  \item \emph{Recall:} say the Bengali for a sigh, then the Bengali for to sob, then say \emph{bilāp kôrā} again
\end{itemize}

\begin{tcolorbox}[breakable,colback=teal!4,colframe=teal!35!black,title={Before you move on}]
What does \bn{বিলাপ} \bn{করা} mean? (\textquotedblleft{}To lament; to wail in grief\textquotedblright{}.) Say it once more.
\end{tcolorbox}
\section[(dialogue)]{Looking back: from Mind and Feeling to Grief and Tears}

\label{lesson:BN-R445-review-438-445}



\begin{quote}
Say the Bengali for an attitude and a feeling.
\end{quote}

\subsection*{Your turn}
Say these aloud:

\begin{itemize}
\raggedright
  \item an attitude, a feeling, emotion, to feel, to become aware of
  \item happiness, suffering, satisfaction, to enjoy, to get carried away
  \item wonder, curiosity, enchanted, to be surprised, to be startled
  \item terror, to recoil in fright, to shudder, to get flustered, to pound
  \item anxiety, anxious thoughts, restless, free of worry, to fidget
  \item resentment, irritable, to provoke, to seethe, to grumble
  \item hurt pride, a gentle complaint, to sulk, to forgive, ill feeling
  \item grief, crying, a sigh, to sob, to lament
\end{itemize}

\begin{tcolorbox}[breakable,colback=teal!4,colframe=teal!35!black,title={Before you move on}]
An attitude? (\textbf{\bn{মনোভাব}}, \emph{mônobhāb}.) A feeling? (\textbf{\bn{অনুভূতি}}, \emph{ônubhūti}.) Emotion? (\textbf{\bn{আবেগ}}, \emph{ābeg}.) To feel? (\textbf{\bn{অনুভব} \bn{করা}}, \emph{ônubhôb kôrā}.) To become aware of? (\textbf{\bn{টের} \bn{পাওয়া}}, \emph{ṭer pāoyā}.) Happiness? (\textbf{\bn{সুখ}}, \emph{sukh}.) Suffering? (\textbf{\bn{কষ্ট}}, \emph{kôshṭô}.) Satisfaction? (\textbf{\bn{তৃপ্তি}}, \emph{tripti}.) To enjoy? (\textbf{\bn{উপভোগ} \bn{করা}}, \emph{upôbhog kôrā}.) To get carried away? (\textbf{\bn{মেতে} \bn{ওঠা}}, \emph{mete oṭhā}.) Wonder? (\textbf{\bn{বিস্ময়}}, \emph{bismôy}.) Curiosity? (\textbf{\bn{কৌতূহল}}, \emph{koutūhôl}.) Enchanted? (\textbf{\bn{মুগ্ধ}}, \emph{mugdhô}.) To be surprised? (\textbf{\bn{অবাক} \bn{হওয়া}}, \emph{ôbāk hôoyā}.) To be startled? (\textbf{\bn{চমকানো}}, \emph{chômkāno}.) Terror? (\textbf{\bn{আতঙ্ক}}, \emph{ātôngkô}.) To recoil in fright? (\textbf{\bn{আঁতকে} \bn{ওঠা}}, \emph{ā̃tke oṭhā}.) To shudder? (\textbf{\bn{শিউরে} \bn{ওঠা}}, \emph{shiure oṭhā}.) To get flustered? (\textbf{\bn{ঘাবড়ানো}}, \emph{ghābṛāno}.) To pound? (\textbf{\bn{ধড়ফড়} \bn{করা}}, \emph{dhôṛphôṛ kôrā}.) Anxiety? (\textbf{\bn{উদ্বেগ}}, \emph{udbeg}.) Anxious thoughts? (\textbf{\bn{দুর্ভাবনা}}, \emph{durbhābônā}.) Restless? (\textbf{\bn{অস্থির}}, \emph{ôsthir}.) Free of worry? (\textbf{\bn{নিশ্চিন্ত}}, \emph{nishchintô}.) To fidget? (\textbf{\bn{ছটফট} \bn{করা}}, \emph{chhôṭphôṭ kôrā}.) Resentment? (\textbf{\bn{ক্ষোভ}}, \emph{khobh}.) Irritable? (\textbf{\bn{খিটখিটে}}, \emph{khiṭkhiṭe}.) To provoke? (\textbf{\bn{খেপানো}}, \emph{khepāno}.) To seethe? (\textbf{\bn{ফুঁসে} \bn{ওঠা}}, \emph{phũse oṭhā}.) To grumble? (\textbf{\bn{গজগজ} \bn{করা}}, \emph{gôjgôj kôrā}.) Hurt pride? (\textbf{\bn{অভিমান}}, \emph{ôbhimān}.) A gentle complaint? (\textbf{\bn{অনুযোগ}}, \emph{ônujog}.) To sulk? (\textbf{\bn{মুখ} \bn{ভার} \bn{করা}}, \emph{mukh bhār kôrā}.) To forgive? (\textbf{\bn{ক্ষমা} \bn{করা}}, \emph{khômā kôrā}.) Ill feeling? (\textbf{\bn{মনোমালিন্য}}, \emph{mônomālinyô}.) Grief? (\textbf{\bn{শোক}}, \emph{shok}.) Crying? (\textbf{\bn{কান্না}}, \emph{kānnā}.) A sigh? (\textbf{\bn{দীর্ঘশ্বাস}}, \emph{dīrghôshvās}.) To sob? (\textbf{\bn{ফোঁপানো}}, \emph{phõpāno}.) To lament? (\textbf{\bn{বিলাপ} \bn{করা}}, \emph{bilāp kôrā}.)
\end{tcolorbox}
